\chapter{The Bus and the Address Map}
\label{ch:bus-and-address-map}

The CPU is the only master of the bus. SIRCADA and SIRCAS are slaves: each answers a bus cycle only when it is
selected, and neither can start one. This chapter gives the signals that make up the bus, how a device is selected, the
address each device has, and how long each one takes to answer. The cycle itself, with its phases and diagrams, belongs
to the \hwmanual{}.

\section{The Signals}
\label{sec:bus-signals}

\begin{table}[H]
	\centering
	\small
	\begin{tabularx}{\textwidth}{llllX}
		\toprule
		\textbf{Signal}              & \textbf{Driven by} & \textbf{Used by}          & \textbf{Type} & \textbf{Meaning}                         \\
		\midrule
		A0--A23                      & CPU                & Select logic, memories    & Three-state   & 24-bit word address                      \\
		A0--A15                      & CPU                & SIRCADA, SIRCAS           & Three-state   & The word offset in the device's segment  \\
		D0--D15                      & CPU or device      & All                       & Three-state   & 16-bit data                              \\
		\activelow{BAS}              & CPU                & Select logic, all devices &               & Marks a bus cycle                        \\
		BRW                          & CPU                & All                       &               & High for a read, low for a write         \\
		\activelow{BACK}             & Devices            & CPU                       & Open drain    & Completes the cycle                      \\
		\activelow{BERR}             & System             & CPU                       &               & Aborts the cycle with a bus fault        \\
		\activelow{BPER}             & System             & CPU                       &               & Aborts the cycle with a protection fault \\
		BAT0--BAT2                   & CPU                & Memory controller, if any &               & The kind of access                       \\
		PROT                         & CPU                & Memory controller, if any &               & High in protected mode                   \\
		\activelow{CS} (each device) & Select logic       & SIRCADA, SIRCAS           &               & Selects the device                       \\
		\bottomrule
	\end{tabularx}
	\caption{The Bus Signals}
	\label{tab:bus-signals}
\end{table}

SIRCADA and SIRCAS take only the sixteen low address bits and have no inputs for \activelow{BERR}, \activelow{BPER},
\texttt{BAT} or \texttt{PROT}. They see the cycle as a chip select, a strobe, a read/write line and an address, and
they answer it with data and \activelow{BACK}.

\section{Selecting a Device}
\label{sec:selecting}

The eight high bits of the address, A16 to A23, are the number of a \emph{segment} of 65,536 words. SIRCUS has no
decoder. Each device is selected by an address bit wired straight to its chip select, so the bit that selects a device
is the one that is set in its segment number. A device is wired to no other address bit above its own 16, so it answers
at every address in which its selecting bit is set, and its segment repeats at every segment above it.

The \activelow{CS} input of SIRCAS is active low, so A23 reaches it through an inverter. SIRCAS then answers whenever
A23 is set, and its first segment is \texttt{0x80}. A device also needs \activelow{BAS} to be asserted: it answers a
cycle only when its \activelow{CS} and \activelow{BAS} are both asserted, and it never drives the data bus at any other
time.

A program uses one segment of each device and does not use an address in which the selecting bits of two devices are
both set: both would answer, and both would drive the data bus. The bits that select the other devices and the memories
are not yet specified.

\section{The Address Map}
\label{sec:address-map}

\begin{table}[H]
	\centering
	\begin{tabular}{llp{6cm}}
		\toprule
		\textbf{Segment}  & \textbf{Device}                & \textbf{Notes}                                                                                                                                           \\
		\midrule
		\texttt{0x00}     & Working memory or system ROM   & The system ROM while the boot overlay is set, the 128 KB of working memory after it                                                                      \\
		Not yet specified & System ROM, at its own segment & Always visible, so its routines can be called after the overlay is cleared                                                                               \\
		Not yet specified & Cartridge                      & The program and data of the game; the connector and the segments are not yet specified                                                                   \\
		\texttt{0x0C}     & SIRCADA                        & Registers, palette, sprite table and video memory (\videomanual{}, memory map chapter)                                                                   \\
		\texttt{0x80}     & SIRCAS                         & Selected by A23, so it also answers in \texttt{0x81} to \texttt{0xFF}; registers and the window to the audio memory (\audiomanual{}, memory map chapter) \\
		Others            & ---                            & Not yet specified                                                                                                                                        \\
		\bottomrule
	\end{tabular}
	\caption{The SIRCUS Address Map}
	\label{tab:address-map}
\end{table}

\subsection{The Vector Table and the Boot Overlay}
\label{sec:boot-overlay}

The CPU fetches its reset vector, and every other exception vector, from a table at words \texttt{0x0000} to
\texttt{0x01FF} of the address space, which cannot be moved (\archmanual{}, exceptions chapter). The same table holds
the vectors of the interrupts and of the traps, among them the handlers for the vertical sync and line-compare
interrupts of the display adapter and for the completion interrupt of SIRCAS. The table must be in ROM when the system
starts, because the reset vector is the first entry. It must be in RAM when a program runs, so that the program can
write its own handlers. A \emph{boot overlay} provides both:

\begin{itemize}
	\item While the overlay is set, the system ROM answers at segment \texttt{0x00}, and the working memory does not.
	\item Reset sets the overlay. \activelow{RSTO} is wired to it, so a software reset sets it too.
	\item A write to the system ROM's own segment clears the overlay. The ROM ignores the write itself. From then on the working
	      memory answers at segment \texttt{0x00}, and the system ROM answers only at its own segment.
\end{itemize}

The overlay is one latch and a few gates. The system ROM's start-up code runs from the low words, jumps to the same
code at the ROM's own segment, clears the overlay by writing there, and writes the default vector table into the
working memory. It then sets the system up and jumps to the program on the cartridge (Section~\ref{sec:system-start}).

From that point the program owns the vector table. It may overwrite the whole table, and a program that wants to do
without the system ROM altogether need never call it again. The ROM stays readable at its own segment for programs that
do use its driver routines. Where the system ROM and the cartridge sit in the address space, and the logic that selects
the memories, are not yet specified.

The contents of the system ROM, its start-up sequence and the interface to its driver routines are not described in
this manual.

\section{What Happens at an Unmapped Address}
\label{sec:unmapped}

A cycle to an address that selects no device is answered by nobody. Which addresses those are depends on the logic that
selects the memories, which is not yet specified. No \activelow{BACK} is asserted, so the CPU inserts wait states and,
unless the system ends the cycle, waits for ever. A system that wants a fault instead must assert \activelow{BERR} for
a segment it has not mapped, which makes the CPU raise a bus fault (\hwmanual{}, bus interface chapter). Whether SIRCUS
does so, and after how long, is not yet specified.

\section{How Long Each Device Takes}
\label{sec:answer-times}

A cycle takes one \texttt{CLKI} cycle if the device asserts \activelow{BACK}, and drives the data of a read, in the
same cycle it is selected. A device that cannot adds wait states, one cycle at a time (\hwmanual{}, bus interface
chapter).

\begin{table}[H]
	\centering
	\begin{tabular}{p{5.4cm}p{2.7cm}p{4.2cm}}
		\toprule
		\textbf{Access}                                        & \textbf{Cycles}   & \textbf{Source}                       \\
		\midrule
		SIRCADA registers, palette, sprite table, video memory & 1, no wait states & \videomanual{}, timing chapter        \\
		SIRCAS registers                                       & 1, no wait states & \audiomanual{}, signal timing chapter \\
		SIRCAS audio memory window                             & 2 to 5            & \audiomanual{}, signal timing chapter \\
		Working memory                                         & Not yet specified &                                       \\
		System ROM                                             & Not yet specified &                                       \\
		Cartridge                                              & Not yet specified &                                       \\
		\bottomrule
	\end{tabular}
	\caption{Bus Cycle Lengths}
	\label{tab:answer-times}
\end{table}

A window access to the audio memory takes more than one cycle because SIRCAS gives the CPU only some of the slots of
the audio memory (\audiomanual{}, timing chapter): the CPU waits until one comes round, and the device asserts
\activelow{BACK} when it has made the access. The wait depends only on where in the slot pattern the request falls,
never on what the voices are doing, so it is at most four wait states.

\section{DMA}

The CPU has a DMA coprocessor that moves blocks of words between addresses with bus cycles of its own (\hwmanual{}, bus
interface chapter). The \texttt{BAT} bits mark them as bursts. SIRCADA and SIRCAS see them as ordinary cycles, one word
at a time, and they answer them the same way. A transfer to the audio window waits two to five cycles for each word it
writes, in addition to the cycle that reads the word from its source. A transfer to the display adapter has no waits: a
transfer of 255 words takes 522 cycles (\videomanual{}, timing chapter).

\section{Protection}

The CPU drives \texttt{PROT} and \texttt{BAT} so that a memory controller can refuse an access, and refuses it by
asserting \activelow{BPER}. SIRCADA and SIRCAS do not use these signals and do not themselves protect any of their
registers. Whether SIRCUS has a memory controller, and which regions it protects, is not yet specified.
